\chapter{Conclusion}
\label{ch:conclusion}

The thesis set out to investigate a graph-based attack-simulation pipeline and automated defense optimizer. The implementation described in Chapter~\ref{ch:implementation} provides a context-graph model, state-transition simulation, and defense strategies. Chapter~\ref{ch:evaluation} defines the evidence required to assess them; it does not report empirical findings.

\section{Summary}
\label{ch:conclusion:summary}

Three contributions were made. First, a \emph{context graph} and state-transition model informed by attack-graph literature, with synthetic CVSS data and independently assigned success probabilities. Second, a \emph{Monte Carlo blast-radius simulator} for the fixed stylized attacker model. Third, a small \emph{optimization framework} supporting policy-rule segmentation, simulation-informed patching, and annealing-based search. Table~\ref{tab:conclusion:contributions} summarizes the contributions and their status.

\begin{table}[hbt!]
	\centering
	\begin{tabular}{lll}
		\toprule
		\textbf{Area}             & \textbf{Contribution}                          & \textbf{Status}      \\
		\midrule
		Context graph             & Informed by attack-graph literature            & Implemented          \\
		Monte Carlo simulation    & Uniform action selection, stylized outcomes    & Implemented          \\
		Optimization framework    & Unit-cost patching, segmentation, and annealing & Implemented         \\
		Empirical evaluation      & Reproducible protocol and reporting requirements & Pending             \\
		\bottomrule
	\end{tabular}
	\caption{Contributions of the thesis and their status.}
	\label{tab:conclusion:contributions}
\end{table}

\section{Future Work}
\label{ch:conclusion:future}

Several directions are left open for future work:

\begin{itemize}
	\item \textbf{Real-CVE integration.} The current vulnerability catalog is synthetic. A live NVD integration~\cite{nvd} would require a separate, justified mapping from CVE data to model parameters; CVSS must not be treated as a calibrated exploit probability.
	\item \textbf{Dynamic topology.} The present system assumes a static network. Extending the attack-graph and simulator layers to handle changes in the network over time would unlock operational use cases (e.g., daily re-evaluation of a corporate network).
	\item \textbf{Richer attacker models.} The current Monte Carlo simulator has one stylized model and no attacker-profile parameter. A more expressive model, perhaps informed by incident data, would require separate validation before it could support claims about attacker behavior~\cite{jha02two, ammann02tva}.
	\item \textbf{Better diagnostics.} Translating the optimizer's recommendations back to source-level positions in the input topology (e.g., \enquote{segment the link between $H_3$ and $H_7$, patch CVE-2024-1234 on $H_{12}$}) would substantially improve the usability of the system for security operators.
\end{itemize}

\section{Personal Growth}
\label{ch:conclusion:growth}

Beyond the technical contributions, the project was a useful exercise in keeping a small, focused design coherent as the surface area grew. The act of writing examples that double as documentation forced several early design decisions to be revisited; in every case the result was a cleaner interface.

\begin{quote}
	\itshape A security tool is only as good as the assumptions it makes explicit. The blast-radius number on its own is not a defense; the act of computing it forces the defender to make the assumptions visible.
\end{quote}
